\section{同一分区上的 KL 粗化与梯度残差}
\label{sec:kl-partition-geometry}
\label{sec:kl-partition}

头部精确求和、尾部合并和尾部残差抽样，都可以从同一个词表分区推导。分区同时规定可精确计算的粗粒度目标和需要抽样的细粒度残差。下面始终固定 prefix、当前参数和参考分布；只在明确写出梯度时对 student 参数求导。

\begin{definition}[词表分区与桶内分布]
令 $\mathcal C$ 是有限词表 $\mathcal V$ 的一个分区，各个单元 $C\in\mathcal C$ 非空且两两不交，并满足 $\bigcup_{C\in\mathcal C}C=\mathcal V$。分区在求导时固定。定义
\[
 P_C=\sum_{v\in C}p(v),\qquad
 Q_C=\sum_{v\in C}q(v),\qquad
 \ell_C=\log\frac{P_C}{Q_C},
\]
以及桶内条件分布和平均 score
\[
 p_C(v)=\frac{p(v)}{P_C},\qquad
 q_C(v)=\frac{q(v)}{Q_C},\qquad
 \bar s_C=\sum_{v\in C}p_C(v)s_v=\nabla\log P_C.
\]
用 $C(v)$ 表示包含 $v$ 的唯一单元。粗粒度反向 KL 与桶内形状 KL 分别定义为
\[
 K_{\mathcal C}=\sum_{C\in\mathcal C}P_C\ell_C,
 \qquad D_C=\KL(p_C\Vert q_C).
\]
这里 $K_{\mathcal C}$ 是桶标签分布之间的 KL；$D_C$ 描述保留桶标签以后仍未解释的条件分布差异。
\end{definition}

\begin{theorem}[同一分区上的双向 KL 链式分解]
\label{thm:kl-partition-chain}
在本章正支持假设下，反向 KL 和正向 KL 分别满足
\begin{align}
 K&=K_{\mathcal C}+\sum_{C\in\mathcal C}P_CD_C,
 \label{eq:kl-partition-chain-reverse}\\
 K^{\mathrm F}&=K^{\mathrm F}_{\mathcal C}
       +\sum_{C\in\mathcal C}Q_CD^{\mathrm F}_C,
 \label{eq:kl-partition-chain-forward}
\end{align}
其中
\[
 K^{\mathrm F}_{\mathcal C}=-\sum_CQ_C\ell_C,
 \qquad D_C^{\mathrm F}=\KL(q_C\Vert p_C).
\]
因此两种方向的粗粒度 KL 都是对应完整 KL 的下界。等号成立当且仅当每个单元内的条件分布相同。
\end{theorem}

\begin{proof}
对任意 $v\in C$，完整 log-ratio 可以逐点写为
\[
 r_v=\log\frac{p(v)}{q(v)}
     =\ell_C+\log\frac{p_C(v)}{q_C(v)}.
\]
先乘以 $p(v)=P_Cp_C(v)$，在单元内求和，再对单元求和，得到
\[
 \sum_vp(v)r_v
 =\sum_CP_C\ell_C
  +\sum_CP_C\sum_{v\in C}p_C(v)
                        \log\frac{p_C(v)}{q_C(v)}.
\]
这就是反向分解。将上面的逐点恒等式取负并乘以 $q(v)=Q_Cq_C(v)$，按相同顺序求和，得到正向分解。所有质量严格为正，且每个条件 KL 非负；条件 KL 等于零恰好对应桶内条件分布相同，故得到下界及等号条件。
\end{proof}

\begin{proposition}[分区细化的标量单调性]
\label{prop:kl-partition-refinement}
设 $\mathcal C'$ 细化 $\mathcal C$。记 $\mathcal C'|_C$ 为包含在 $C$ 内的细单元集合，则
\[
 K_{\mathcal C'}-K_{\mathcal C}
 =\sum_{C\in\mathcal C}P_C
   \KL\left(
       \left(\frac{P_{C'}}{P_C}\right)_{C'\in\mathcal C'|_C}
       \middle\Vert
       \left(\frac{Q_{C'}}{Q_C}\right)_{C'\in\mathcal C'|_C}
      \right)\geq0.
\]
正向版本交换两组质量，并将外层权重改为 $Q_C$。因此细化使粗粒度 KL 的标量偏差绝对值和平方偏差单调下降；全部单元都是 singleton 时偏差为零。
\end{proposition}

\begin{proof}
对每个 $C$，利用 $\sum_{C'\subseteq C}P_{C'}=P_C$，将新增贡献写成
\begin{align*}
 \sum_{C'\subseteq C}P_{C'}\log\frac{P_{C'}}{Q_{C'}}
    -P_C\log\frac{P_C}{Q_C}
 &=P_C\sum_{C'\subseteq C}\frac{P_{C'}}{P_C}
       \log\frac{P_{C'}/P_C}{Q_{C'}/Q_C}.
\end{align*}
右边是 $P_C$ 乘以细单元标签之间的条件 KL。对 $C$ 求和即得所述等式。结合粗粒度下界，误差 $K-K_{\mathcal C}$ 非负且随细化下降，其平方也下降。
\end{proof}

\begin{theorem}[粗粒度 KL 的真实梯度及其残差]
\label{thm:kl-partition-gradient}
记 $g_{\mathcal C}=\nabla K_{\mathcal C}$。则
\begin{align}
 g_{\mathcal C}
 &=\sum_CP_C\ell_C\bar s_C
  =\sum_vp(v)\ell_{C(v)}s_v,
 \label{eq:kl-partition-gradient}\\
 g-g_{\mathcal C}
 &=\sum_CP_C\E_{p_C}[(r_A-\ell_C)s_A]
 \label{eq:kl-partition-gradient-residual}\\
 &=\sum_CP_C\left\{
       \operatorname{Cov}_{p_C}(r_A-\ell_C,s_A)
       +D_C\bar s_C\right\}.
 \label{eq:kl-partition-gradient-covariance}
\end{align}
标量与向量的协方差定义为
\[
 \operatorname{Cov}_{p_C}(u_A,s_A)
 =\E_{p_C}[(u_A-\E_{p_C}u_A)(s_A-\bar s_C)].
\]
\end{theorem}

\begin{proof}
由于 $q$ 固定，$\nabla Q_C=0$，而 $\nabla P_C=P_C\bar s_C$。逐项求导得到
\begin{align*}
 \nabla K_{\mathcal C}
 &=\sum_C\{\ell_C\nabla P_C+P_C\nabla\ell_C\}\\
 &=\sum_C\ell_C\nabla P_C+\sum_C\nabla P_C\\
 &=\sum_CP_C\ell_C\bar s_C.
\end{align*}
最后一步用了 $\sum_C\nabla P_C=\nabla1=0$。将 $P_C\bar s_C$ 展开为桶内加权 score，得到第二种表示。完整梯度为 $g=\sum_vp(v)r_vs_v$，与粗粒度梯度逐项相减，得到式~\eqref{eq:kl-partition-gradient-residual}。又因为
\[
 \E_{p_C}[r_A-\ell_C]=D_C,
\]
用 $\E[us]=\operatorname{Cov}(u,s)+\E[u]\E[s]$ 展开每个桶内期望，就得到式~\eqref{eq:kl-partition-gradient-covariance}。
\end{proof}

式~\eqref{eq:kl-partition-gradient-covariance} 把梯度误差分为桶内相关项和条件形状 KL 乘以平均 score 的项。标量粗化误差只包含 $D_C$ 的加权和，梯度误差还依赖 score 的方向与桶内相关性。因此标量下界的单调改进不能直接推出梯度偏差范数单调下降；梯度结论需要单独分析这些向量项。

对正向 KL，冻结权重 $Q_C$ 使粗粒度梯度变成
\[
 g^{\mathrm F}_{\mathcal C}=-\sum_CQ_C\bar s_C,
 \qquad
 g^{\mathrm F}-g^{\mathrm F}_{\mathcal C}
 =-\sum_CQ_C\{\E_{q_C}[s_A]-\bar s_C\}.
\]
这里平均 student score $\bar s_C$ 的积分分布仍是 $p_C$；它不能替换成 teacher 条件均值。第一式来自 $\nabla\ell_C=\bar s_C$，第二式来自完整正向梯度 $-\sum_vq(v)s_v$。

\subsection{一个标量代理与两个同积分的梯度代理}

使用前文的精确积分加残差抽样构造，令标量代理为
\[
 a^{\mathrm K}_v=\ell_{C(v)},\qquad
 I_p(a^{\mathrm K})=K_{\mathcal C}.
\]
在 $A\sim p$ 下，相应无偏数值估计量为
\[
 \widehat K_{\mathcal C}
 =K_{\mathcal C}+r_A-\ell_{C(A)}.
\]
它的均值由链式分解得到完整 $K$，方差为 $\Var_p(r_A-\ell_{C(A)})$。桶内条件形状相同时，残差逐点为零，估计量成为确定性的完整 KL。

梯度积分的目标函数为 $r_vs_v$。同一个粗粒度梯度允许两种可精确积分的代理：
\begin{align*}
 a^{\mathrm{G,const}}_v&=\ell_{C(v)}\bar s_{C(v)},\\
 a^{\mathrm{G,score}}_v&=\ell_{C(v)}s_v.
\end{align*}
两者逐点不同，但精确积分相同：
\[
 I_p(a^{\mathrm{G,const}})
 =\sum_CP_C\ell_C\bar s_C
 =I_p(a^{\mathrm{G,score}})=g_{\mathcal C}.
\]
因此它们产生两个不同的无偏梯度估计量：
\begin{align}
 \widehat g_{\mathcal C}^{\mathrm{const}}
 &=g_{\mathcal C}+r_As_A-\ell_{C(A)}\bar s_{C(A)},
 \label{eq:kl-partition-constant-estimator}\\
 \widehat g_{\mathcal C}^{\mathrm{score}}
 &=g_{\mathcal C}+(r_A-\ell_{C(A)})s_A.
 \label{eq:kl-partition-score-estimator}
\end{align}
两种代理之差为
\[
 \psi_v=\ell_{C(v)}(s_v-\bar s_{C(v)}),
 \qquad \E_{p_C}[\psi_A]=0.
\]
所以改变这两个代理之间的线性组合，只改变桶内零均值波动。它不改变粗粒度精确积分，也不改变每个桶内残差的条件均值。在前文控制矩阵与 Gram 二次型框架中，$\psi_A$ 就是一列已知均值为零的控制特征；其最优系数由残差与该特征的协方差决定，取系数一不自动保证降低方差。

\subsection{头部 singleton 与尾部单桶的特例}

令 $S$ 为冻结的 teacher Top-$J$ 支持，定义 $P=\sum_{v\in S}p(v)$、$Q=\sum_{v\in S}q(v)$、$\tau_p=1-P$、$\tau_q=1-Q$ 和 $J_A=\mathbf1\{A\notin S\}$，本节头尾公式假设 $0<P,Q<1$。头部积分为
\[
 H=\sum_{v\in S}p(v)r_v,\qquad h_g=\sum_{v\in S}p(v)r_vs_v.
\]
令 $\mathcal C=\{\{v\}:v\in S\}\cup\{\mathcal V\setminus S\}$。头部单元满足 $\ell_{\{v\}}=r_v$、$\bar s_{\{v\}}=s_v$；尾部单元满足
\[
 \ell_{\tail}=\log\frac{\tau_p}{\tau_q},\qquad
 \bar s_\tail=\nabla\log\tau_p,\qquad
 b_{\mathrm{TA}}=\ell_{\tail}\bar s_\tail.
\]
因此 $K_{\mathcal C}=H+\tau_p\ell_{\tail}$ 是 Tail-aware KL，且 $g_{\mathcal C}=h_g+\tau_pb_{\mathrm{TA}}$。三个残差构造分别成为
\begin{align*}
 \widehat K_{\mathcal C}
 &=H+\tau_p\ell_{\tail}+J_A(r_A-\ell_{\tail}),\\
 \widehat g_{\mathcal C}^{\mathrm{const}}
 &=h_g+\tau_pb_{\mathrm{TA}}+J_A(r_As_A-b_{\mathrm{TA}}),\\
 \widehat g_{\mathcal C}^{\mathrm{score}}
 &=h_g+\tau_pb_{\mathrm{TA}}+J_A(r_A-\ell_{\tail})s_A.
\end{align*}
第一式是标量残差估计量；第二式是 ResOPD；第三式是 SC-TA。HETS 使用头部精确代理和零尾部代理，因而是另一个代理选择。上述分类从积分目标和代理函数同时确定估计量，不依赖算法名称。

\begin{center}
\begin{tabular}{lll}
\toprule
目标与方法 & 头部近似 $a_v$ & 尾部近似 $a_v$\\
\midrule
数值 HETS & $r_v$ & $0$\\
数值桶残差校正 & $r_v$ & $\ell_{\tail}$\\
梯度 HETS & $r_vs_v$ & $0$\\
梯度 ResOPD & $r_vs_v$ & $b_{\mathrm{TA}}$\\
梯度 SC-TA & $r_vs_v$ & $\ell_{\tail}s_v$\\
\bottomrule
\end{tabular}
\end{center}
表中每一行都使用式~\eqref{eq:kl-master-estimator}；确定积分项是该行的 $I_p(a)$，采样项始终是 $f_A-a_A$。数值 Tail-aware 与梯度粗化则只返回各自行的确定积分项，没有补回残差。

\begin{proposition}[标量精确时，常向量代理仍可有方差]
若 $p_\tail=q_\tail$，则上述标量残差估计量和 SC-TA 都是确定性的完整目标，而 ResOPD 的协方差为
\[
 \operatorname{Cov}(\widehat g^{\mathrm{ResOPD}})
 =\tau_p\ell_{\tail}^2\operatorname{Cov}_{p_\tail}(s_A).
\]
\end{proposition}

\begin{proof}
桶内条件分布相同意味着尾部 $r_v=\ell_{\tail}$。标量残差和 SC-TA 的随机残差于是逐点为零。ResOPD 则仍有随机残差
\[
 J_A\ell_{\tail}(s_A-\bar s_\tail).
\]
它在尾部条件下均值为零，在头部为零，因此总体均值为零，二阶矩恰好是命题中的协方差。
\end{proof}

例如取 $p=(0.4,0.2,0.4)$、$q=(0.7,0.1,0.2)$ 和 $S=\{1\}$，则两个尾部条件分布均为 $(1/3,2/3)$，且 $\tau_p=0.6$、$\ell_{\tail}=\log2$。在完整 softmax logits 参数化下，$s_v=e_v-p$，尾部 score 的协方差迹为 $1-(1/3)^2-(2/3)^2=4/9$。ResOPD 的方差因此为 $\frac4{15}(\log2)^2>0$；标量残差估计量与 SC-TA 的方差均为零。这个例子直接区分了可精确积分的 KL 数值、粗粒度真实梯度以及不同梯度代理留下的抽样波动。
